\title{xLSTM: Extended Long Short-Term Memory}

\begin{abstract}
The 1990s saw the introduction of the constant error carousel and gating mechanisms, which became foundational components of the Long Short-Term Memory (LSTM) architecture. LSTMs have remained relevant over time and have played a pivotal role in many major advances in deep learning, notably serving as the backbone of the earliest Large Language Models (LLMs). Nevertheless, the rise of Transformer models—with their scalable and parallelizable self-attention mechanisms—ushered in a transformative shift, surpassing LSTMs in large-scale applications. In this work, we pose a central inquiry: to what extent can LSTM-based models excel in language modeling if expanded to the billion-parameter regime, adapted with innovations from contemporary LLM development, and revised to circumvent the traditional limitations of LSTMs? To this end, we present exponential gating alongside normalization and stabilization strategies tailored for large models. Furthermore, we alter the canonical LSTM memory scheme to create: (i) sLSTM, which employs a scalar-valued memory, scalar updates, and revised memory integration, and (ii) mLSTM, which enables full parallelization through matrix-based memories and a covariance-inspired memory update mechanism. By incorporating these enhanced LSTM variants into residual backbone structures, we define xLSTM blocks, which are composed in a stacked residual fashion to build complete xLSTM architectures. These advancements—exponential gating and reengineered memory configurations—allow xLSTM to achieve competitive or superior results relative to current state-of-the-art Transformer models and State Space Models, in both empirical performance and scalability.\\
Code available at: \url{https://github.com/NX-AI/xlstm}
\end{abstract}

\section{Introduction}
\label{sec:intro}

The concepts behind Long Short-Term Memory (LSTM) networks~\citep{Hochreiter:91,Hochreiter:97nips,Hochreiter:97}, namely the use of a constant error carousel and gating mechanisms, were developed as solutions to the vanishing gradient issue prevalent in recurrent neural networks~\citep{Hochreiter:91,Hochreiter:00book}:

\begin{align}
\label{eq:lstm_idea}
  c_t \ = \ f_t \ c_{t-1} \ + \ 
          i_t \  z_t \ , \quad  
          h_t \ = \ o_t \ \psi({c_t}) \ . 
\end{align}

The constant error carousel involves the cell state $c_{t-1}$ (green) being updated additively with the cell inputs $z_t$, with this process regulated by the sigmoid gates (blue). The forget gate $f_t$ and input gate $i_t$ determine the nature of the cell state update, whereas the output gate $o_t$ dictates how information is emitted from the memory cell, producing the hidden state $h_t$. Before being relayed, the cell state is subjected to normalization or a nonlinear transformation by $\psi$, after which the output gate determines the resultant hidden state.

LSTMs have found extensive utility across multiple domains \citep{Hochreiter:01,Hochreiter:07,Schmidhuber:15}, and maintained dominance in text generation until Transformers emerged in 2017~\citep{Vaswani:17}. Their capacity has been proven on a wide array of sequence-based tasks, including text generation~\citep{Graves:13,Karpathy:15blog}, handwriting synthesis~\citep{Graves:13}, sequence-to-sequence translation~\citep{Sutskever:14nips}, program evaluation~\citep{Zaremba:14arxiv}, image captioning~\citep{Karpathy:15,Hossain:19}, code generation~\citep{Karpathy:15blog}, rainfall-runoff prediction~\citep{Kratzert:18,Kratzert:19}, and hydrological flood warning modeling~\citep{Nearing:24}. Within reinforcement learning, LSTMs represent the state of the art among sequence models, exemplified by their integration into AlphaStar for StarCraft II~\citep{Vinyals:17short}, OpenAI Five for Dota 2~\citep{OpenAI:19}, and control models in nuclear fusion magnetic systems~\citep{Degrave:22short}. A notable strength of LSTMs lies in their capacity for abstraction—they efficiently learn and internalize semantic content within their cell states~\citep{Karpathy:15blog}, as shown by discoveries such as number and syntax neurons~\citep{Lakretz:19}, neurons representing linguistic features~\citep{Bau:19}, and neurons relating to sentiment~\citep{Radford:17arxiv}. LSTMs continue to be leveraged for important, contemporary applications~\citep{Degrave:22short,Nearing:24}, reflecting their enduring relevance and robustness.

Despite their notable achievements, LSTMs face three principal drawbacks: (i) Inability to modify prior storage decisions. This shortcoming is illustrated by the {\it Nearest Neighbor Search} task (refer to Appendix~\ref{sec:appNearestNeighborSearch}): Given a reference vector, the model must traverse a sequence to identify and recall the most similar vector and its associated value at the sequence’s conclusion. The left panel of Figure~\ref{fig:lstmProblems} depicts the mean squared error observed for this task. Standard LSTM architectures struggle to replace previously retained values when subsequently encountering vectors of greater similarity, whereas the proposed xLSTM architecture addresses this through exponential gating mechanisms. (ii) Constrained storage abilities, specifically because relevant information is restricted to scalar cell states. This is evident in {\it Rare Token Prediction}. As shown in the right panel of Figure~\ref{fig:lstmProblems}, the perplexity of token prediction on the Wikitext-103 dataset~\citep{Merity:17} is evaluated across token frequency partitions. LSTM demonstrates inferior performance on infrequent tokens owing to its storage bottleneck. Our xLSTM architecture overcomes this by employing a matrix-based memory structure. (iii) Inherent lack of parallelism resulting from memory mixing—that is, the hidden-to-hidden state dependencies between successive time steps—which necessitate sequential data processing.

The shortcomings of LSTM architectures have contributed to the rise of Transformers~\citep{Vaswani:17} within the domain of language modeling. An intriguing question thus arises: if these LSTM limitations are addressed and LSTMs are scaled up to match the parameter sizes typical of contemporary Large Language Models, what level of performance can they reach in language modeling tasks?

\section{Extended Long Short-Term Memory}
\label{sec:xLSTM}

To address the limitations inherent to LSTM architectures, Extended Long Short-Term Memory (xLSTM) proposes two principal advancements to the foundational LSTM concept described in Equation~\eqref{eq:lstm_idea}. These enhancements---namely exponential gating and the adoption of advanced memory structures---expand the LSTM family, introducing (i) sLSTM (refer to Section~\ref{sec:sLSTM}), which is characterized by the use of a scalar memory, scalar updates, and the mechanism of memory mixing, and (ii) mLSTM (see Section~\ref{sec:mLSTM}), which incorporates a matrix-based memory and applies covariance (outer product) update rules, offering full parallelizability. Both sLSTM and mLSTM leverage exponential gating for improved performance. The design of mLSTM forgoes memory mixing---specifically the hidden-to-hidden recurrent pathways---to facilitate parallel computations. The architectures of both sLSTM and mLSTM are amenable to extensions involving multiple memory cells; in the case of sLSTM, this entails memory mixing between cells. Additionally, sLSTM is capable of supporting multiple heads, where memory mixing is restricted to occur among cells within an individual head, and not between heads. Introducing heads to sLSTM, combined with exponential gating, creates a new paradigm for memory mixing. In mLSTM, the notions of multiple heads and multiple cells are effectively interchangeable.

The incorporation of these novel LSTM variants within residual block modules yields what we refer to as xLSTM blocks (refer to Section~\ref{sec:xLSTMblock}). By sequentially layering such xLSTM blocks using residual connections, we construct what we term xLSTM architectures (as discussed in Section~\ref{sec:xLSTMarchitecure}). Figure~\ref{fig:xlstm_sketch} illustrates the xLSTM architecture along with its constituent elements.

\subsection{Review of the Long Short-Term Memory}
\label{sec:orgLSTM}

The foundational concept of LSTM~\citep{Hochreiter:91,Hochreiter:97nips,Hochreiter:97} proposed the scalar memory cell as the core element for processing and storing information, effectively mitigating the issue of vanishing gradients~\citep{Hochreiter:91,Hochreiter:00book} via the mechanism known as the constant error carousel (i.e., cell state updates). Within the memory cell architecture, three types of gates are defined: the input gate, output gate, and forget gate. The introduction of the forget gate is attributed to \citet{Gers:00}. The governing equations that describe the LSTM memory cell’s update process at each time step $t$ are:

\begin{align}
c_t \ &= \  f_t \ c_{t-1} \ + \ i_t \ z_t &  & &\text{cell state} \\
h_t  \ &= \ o_t \ \tilde{h}_t \ , & \tilde{h}_t \ &= \ \psi \left( c_t\right)
  &\text{hidden state} \\
z_t \ &= \ \varphi \left( \tilde{z}_t \right) \ , 
  &\tilde{z}_t \ &=  \ \mathbf{w}^\top_{z} \ \mathbf{x}_t \ + \
  r_{z}  h_{t-1} \ + \  b_{z} \ \
  &\text{cell input} \\
i_t \ &= \ \sigma \left( \tilde{i}_t  \right) \ , 
  &\tilde{i}_t \ &= \ \mathbf{w}^\top_{i} \ \mathbf{x}_t \ + \
  r_{i}  \ h_{t-1} \ + \  b_{i} \ \
  &\text{input gate} \\
f_t \ &= \ \sigma \left(  \tilde{f}_t \right) \ , 
  &\tilde{f}_t \ &= \ \mathbf{w}^\top_{f} \ \mathbf{x}_t  \ + \
  r_{f}  \ h_{t-1} \ + \  b_{f} \ \
  &\text{forget gate} \\
o_t \ &= \ \sigma \left( \tilde{o}_t \right) \ , 
  &\tilde{o}_t  \ &= \ \mathbf{w}^\top_{o} \ \mathbf{x}_t \ + \
  r_{o}  \ h_{t-1} \ + \  b_{o} \ \
  &\text{output gate} 
\end{align}

The vectors $\mathbf{w}_{z}$, $\mathbf{w}_{i}$, $\mathbf{w}_{f}$, and $\mathbf{w}_{o}$ denote the input weight parameters connecting the input $\mathbf{x}_t$ to the cell input, input gate, forget gate, and output gate, in that order. The parameters $r_{z}$, $r_{i}$, $r_{f}$, and $r_{o}$ represent the recurrent weights from the previous hidden state $h_{t-1}$ to each respective element: cell input, input gate, forget gate, and output gate. The terms $b_{z}$, $b_{i}$, $b_{f}$, and $b_{o}$ serve as the biases for the associated gates and cell input. The functions $\varphi$ and $\psi$ are non-linearities applied to the cell input and the hidden state—commonly $\tanh$—with $\psi$ serving to regulate or bound the otherwise unconstrained cell state. All gate nonlinearities use the sigmoid activation function defined by $\sigma \left( x \right)= 1/(1+\exp(-x))$. Later developments introduced the aggregation of several scalar memory cells $\mathbf{c}_t \in \mathbb{R}^{d}$ into a single vector, enabling the application of recurrent matrices $\mathbf{R} \in \mathbb{R}^{d \times d}$, which facilitates interaction among the outputs of the various memory cells~\citep{Greff:15}; additional discussion is provided in Appendix~\ref{sec:apporgLSTM}. Empirical ablation analyses have demonstrated that every constituent component of the memory cell plays an essential role~\citep{Greff:15}.

\subsection{sLSTM}
\label{sec:sLSTM}

In order to enable LSTMs to modify their storage choices, we incorporate exponential gates (highlighted in red) along with mechanisms for normalization and stabilization. Specifically, the input and forget gates may employ exponential activation functions. To achieve normalization, we propose a normalizer state, which accumulates the sum of the input gate multiplied by the sequence of subsequent forget gates.

The scalar sLSTM forward pass is:

\begin{align}
\label{eq:slstmforward}
c_t \ &= \  f_t \ c_{t-1} \ + \ i_t \ z_t & & 
 &\text{cell state} \\
n_t \ &= \  f_t \ n_{t-1} \ + \ i_t  & & 
 &\text{normalizer state} \\
h_t \ &= \ o_t \ \tilde{h}_t \ ,
  &\tilde{h}_t \ &= \ c_t / n_t
&\text{hidden state} \\
z_t \ &= \ \varphi \left( \tilde{z}_t \right) \ , 
  &\tilde{z}_t \ &=  \ \mathbf{w}^\top_{z} \ \mathbf{x}_t \ + \
  r_{z}  h_{t-1} \ + \  b_{z}
  &\text{cell input} \\
i_t \ &= \ \!\exp \left( \tilde{i}_t  \right) \ , 
  &\tilde{i}_t \ &= \ \mathbf{w}^\top_{i} \ \mathbf{x}_t \ + \
  r_{i}  \ h_{t-1} \ + \  b_{i}
  &\text{input gate} \\
f_t \ &= \ \sigma \left(  \tilde{f}_t \right) \ \text{OR} \
\!\exp \left(  \tilde{f}_t \right) \ ,
  &\tilde{f}_t \ &= \ \mathbf{w}^\top_{f} \ \mathbf{x}_t  \ + \
  r_{f}  \ h_{t-1} \ + \  b_{f}
  &\text{forget gate} \\
o_t \ &= \ \sigma \left( \tilde{o}_t \right) \ , 
  &\tilde{o}_t  \ &= \ \mathbf{w}^\top_{o} \ \mathbf{x}_t \ + \
  r_{o}  \ h_{t-1} \ + \  b_{o}
  &\text{output gate} 
\end{align}

---

\begin{align}
\label{eq:slstmforward}
c_t \ &= \  f_t \ c_{t-1} \ + \ i_t \ z_t & & 
 &\text{cell state update} \\
n_t \ &= \  f_t \ n_{t-1} \ + \ i_t  & & 
 &\text{normalization factor update} \\
h_t \ &= \ o_t \ \tilde{h}_t \ ,
  &\tilde{h}_t \ &= \ c_t / n_t
&\text{computation of hidden state} \\
z_t \ &= \ \varphi \left( \tilde{z}_t \right) \ , 
  &\tilde{z}_t \ &=  \ \mathbf{w}^\top_{z} \ \mathbf{x}_t \ + \
  r_{z}  h_{t-1} \ + \  b_{z}
  &\text{candidate cell content} \\
i_t \ &= \ \!\exp \left( \tilde{i}_t  \right) \ , 
  &\tilde{i}_t \ &= \ \mathbf{w}^\top_{i} \ \mathbf{x}_t \ + \
  r_{i}  \ h_{t-1} \ + \  b_{i}
  &\text{input gate value} \\
f_t \ &= \ \sigma \left(  \tilde{f}_t \right) \ \text{OR} \
\!\exp \left(  \tilde{f}_t \right) \ ,
  &\tilde{f}_t \ &= \ \mathbf{w}^\top_{f} \ \mathbf{x}_t  \ + \
  r_{f}  \ h_{t-1} \ + \  b_{f}
  &\text{forget gate value} \\
o_t \ &= \ \sigma \left( \tilde{o}_t \right) \ , 
  &\tilde{o}_t  \ &= \ \mathbf{w}^\top_{o} \ \mathbf{x}_t \ + \
  r_{o}  \ h_{t-1} \ + \  b_{o}
  &\text{output gate value} 
\end{

The established LSTM gating mechanisms, which involve input- and/or hidden-state-dependent gating along with a bias component, are adapted for the novel architectures. Since exponential activations have the potential to generate excessively large outputs resulting in overflows, we introduce an extra state \mbox{$m_t$~\citep{Milakov:18arxiv}} to enhance gate stability:

\begin{align}
\label{eq:slstmstabil}
m_t \ &= \ \max \bigl( \log ( f_t ) + m_{t-1} , \log ( i_t ) \bigr) &\text{stabilizer state} \\
i'_t \ &= \ \!\exp \left( \log \bigl( i_t \bigr) - m_t \right) \ = \ \!\exp \left( \tilde{i}_t - m_t \right) \ \, 
  &\text{stabil. input gate} \\
f'_t \ &= \ \!\exp \left( \log \bigl( f_t \bigr) + m_{t-1} - m_t \right) \ \, 
  &\text{stabil. forget gate}
\end{align}

As demonstrated in Appendix~\ref{sec:appsLSTM}, substituting $f_t$ with $f'_t$ and $i_t$ with $i'_t$ during the forward pass leaves both the network's overall output and the gradients of the loss with respect to the parameters unaffected.

\paragraph{New Memory Mixing.} Similar to the standard LSTM, sLSTM is capable of utilizing several memory cells (refer to Appendix~\ref{sec:appsLSTM}). The presence of multiple memory cells facilitates memory mixing through the recurrent matrices $\mathbf{R}_{\mathbf{z}}$, $\mathbf{R}_{\mathbf{i}}$, $\mathbf{R}_{\mathbf{f}}$, and $\mathbf{R}_{\mathbf{o}}$, which connect the hidden state vector $\mathbf{h}$ to the memory cell input $\mathbf{z}$ as well as the gates $\mathbf{i}$, $\mathbf{f}$, and $\mathbf{o}$. A distinguishing characteristic in the context of memory mixing is the inclusion of exponential gating. Within each head, sLSTM incorporates memory mixing, though no such mixing occurs between different heads. The combination of the head structure in sLSTM and the use of exponential gating yields a novel approach to mixing memory.

\subsection{mLSTM}
\label{sec:mLSTM}

To improve the capacity of LSTM architectures, we replace the traditional scalar memory cell $c \in \mathbb{R}$ with a matrix representation $\mathbf{C} \in \mathbb{R}^{d \times d}$. Consequently, information retrieval relies on matrix multiplication. At time step $t$, a tuple consisting of a key $\mathbf{k}_t \in \mathbb{R}^d$ and a value $\mathbf{v}_t \in \mathbb{R}^d$ (adopting nomenclature from Transformer models) is stored. Later, at time $t + \tau$, a query vector $\mathbf{q}_{t+\tau} \in \mathbb{R}^d$ is utilized to recover the value $\mathbf{v}_t$. This approach aligns with the framework introduced in Bidirectional Associative Memories (BAMs) \citep{Kohonen:72,Anderson:72,Nakano:72,Anderson:77}.

The covariance update rule~\citep{Sejnowski:77,Dayan:91} for storing a key-value pair is

\begin{align}
\mathbf{C}_t \ &= \ \mathbf{C}_{t-1} \ + \ \mathbf{v}_{t} \ \mathbf{k}_t^\top \ .
\end{align}

We posit that input normalization via layer-norm precedes their projection onto key and value spaces, ensuring these vectors exhibit zero mean. The covariance update procedure achieves optimality~\citep{Dayan:91} in terms of maximizing the distinctness among retrieved binary representations, which translates to an optimal signal-to-noise ratio. Greater separability can be achieved when retrieval is restricted to pairwise interactions, though this requires accepting the quadratic computational overhead associated with attention mechanisms~\citep{Krotov:16,Krotov:17,Ramsauer:21}. This update procedure aligns with the methodology used in Fast Weight Programmers~\citep{Schmidhuber:92ncfastweights,Schlag:21}, which were subsequently modified to incorporate a fixed decay factor applied to $\mathbf{C}_{t-1}$, along with a fixed learning rate scaling 
$\mathbf{v}_{t} \mathbf{k}_t ^\top$~\citep{Ba:16}.
Following this line of work, we embed the covariance-based update into the LSTM architecture; here, the decay parameter is realized through the forget gate, the learning rate via the input gate, and the output gate governs the scaling of the vector being retrieved.

In this matrix memory model, the normalizer state represents a sum of key vectors, each scaled by both the input gate and the cumulative product of subsequent forget gates. This setup allows the normalizer state to capture the gating strengths over time. Because the dot product between the query and the normalizer state might approach zero, we utilize the absolute value of this dot product and enforce a minimum threshold (commonly 1.0), following prior work~\citep{Sun:23arxiv}. The mLSTM forward pass is as follows:

\begin{align}
\label{eq:mlstm_recurrent_begin}
\mathbf{C}_t \ &= \ f_t \ \mathbf{C}_{t-1} \ + \ 
  i_t \ \mathbf{v}_t \ \mathbf{k}_t^\top & &  &\text{cell state} \\
\mathbf{n}_t \ &= \ f_t \ \mathbf{n}_{t-1} \ + \ 
  i_t \ \mathbf{k}_t & &  &\text{normalizer state} \\
\mathbf{h}_t \ &= \ \mathbf{o}_t \odot \tilde{\mathbf{h}}_t
\ , \qquad \qquad 
\tilde{\mathbf{h}}_t = \frac{\mathbf{C}_t \mathbf{q}_t}{\max \left\{ |\mathbf{n}_t^\top \mathbf{q}_t|, 1 \right\}} 
   &&&\text{hidden state} \\
\mathbf{q}_t \ &= \ \mathbf{W}_q \mathbf{x}_t + \mathbf{b}_q & & &\text{query input} \\
\mathbf{k}_t \ &= \ \frac{1}{\sqrt{d}} \mathbf{W}_k \mathbf{x}_t + \mathbf{b}_k  & & &\text{key input} \\
\mathbf{v}_t \ &= \ \mathbf{W}_v \mathbf{x}_t + \mathbf{b}_v & & &\text{value input} \\
i_t \ &= \ \!\exp\! \left( \tilde{i}_t \right), 
 \qquad \qquad 
 \tilde{i}_t \ = \ \mathbf{w}_i^\top \mathbf{x}_t + b_i
  &&&\text{input gate} \\
f_t \ &= \ \sigma( \tilde{f}_t )\ \text{OR}\ \exp(\tilde{f}_t ), \ \ \
\tilde{f}_t \ = \ \mathbf{w}_f^\top \mathbf{x}_t + b_f
  &&&\text{forget gate} \\
\label{eq:mlstm_recurrent_end}
\mathbf{o}_t \ &= \ \sigma(\tilde{\mathbf{o}}_t) , \qquad \qquad 
\tilde{\mathbf{o}}_t = \mathbf{W}_{\mathbf{o}} \mathbf{x}_t + \mathbf{b}_{\mathbf{o}}
  &&&\text{output gate}
\end{align}

The architecture of mLSTM allows for several memory cells, similar to what is possible with the standard LSTM model. In mLSTM, the presence of multiple heads functions identically to multiple memory cells because no mixing of memory states occurs. To address the instability that can arise from the exponential gating mechanism in mLSTM, we apply stabilization strategies analogous to those utilized in sLSTM (see Equation~\ref{eq:slstmstabil}). Given the absence of memory mixing within mLSTM, its recurrence relations can be restructured to allow for parallel computation. Additional information is provided in Appendix~\ref{sec:appmLSTM}.

\subsection{xLSTM Architecture}
\label{sec:xLSTMarch}

\paragraph{xLSTM Blocks.}
\label{sec:xLSTMblock}
The primary function of an xLSTM block is to encode historical information in a non-linear fashion within a high-dimensional representation, thereby enhancing the model’s ability to distinguish between different contexts or sequences. Distinguishing among various histories is fundamental for accurately forecasting subsequent elements, such as upcoming tokens. In support of this, we refer to Cover’s Theorem~\citep{Cover:65}, which posits that patterns that are mapped into a higher-dimensional space through non-linear transformations are more readily separable by linear decision boundaries compared to the original lower-dimensional space. 

We investigate two designs for residual block architectures: (i) The first utilizes a post up-projection residual block, akin to that used in Transformers. Here, the sequence’s past is first encoded non-linearly within the original dimensionality, then projected linearly to a higher-dimensional space, followed by the application of a non-linear activation and a subsequent linear down-projection back to the original space; this can be seen in the leftmost panel of Figure~\ref{fig:Backbones} and the third column of Figure~\ref{fig:xlstm_sketch}. A more granular illustration appears in Figure~\ref{fig:appsLSTM_detailed} in the appendix. (ii) The second architecture is a residual block that employs a pre up-projection, following the approach used in State Space Models, whereby information is first projected linearly to a higher-dimensional space,

The approach first applies a nonlinear compression of historical information into a high-dimensional representation, which is subsequently transformed back to the initial space using a linear projection. When the xLSTM block is built with an sLSTM, the up-projection step is performed after this compression. Conversely, for xLSTM blocks that utilize an mLSTM, the up-projection precedes the main operation due to the increased memory available in the high-dimensional context. Further clarification can be found by consulting the left side of Figure~\ref{fig:Backbones}, the third column of Figure~\ref{fig:xlstm_sketch}, or Figure~\ref{fig:appsLSTM_detailed} within the appendix.

\paragraph{xLSTM Architecture. }
\label{sec:xLSTMarchitecure}
The xLSTM model is designed by sequentially arranging modular components in a residual stacking manner~\citep{Srivastava:15,He:16}. Our implementation adopts the widely-used pre-LayerNorm~\citep{Ba:16b} residual backbone architecture, consistent with the designs of modern Large Language Models. Refer to the final column of Figure~\ref{fig:xlstm_sketch} for a visual overview.

\subsection{Memory and Speed Considerations}

Unlike Transformers, xLSTM networks feature linear computational complexity and maintain constant memory requirements as sequence length increases. Owing to the compressive nature of xLSTM memory, these networks are particularly advantageous for industrial usage and edge device deployment.

Although the mLSTM's memory does not involve additional parameters, it incurs significant computational costs due to its $d \times d$ matrix structure and the associated $d \times d$ updates. This architecture prioritizes increased memory capacity at the expense of heightened computational demands. However, since these operations can be efficiently parallelized on GPUs, the impact on actual execution time is minimal.

Although mLSTM, similar to FlashAttention~\citep{Dao:22,Dao:23} or GLA~\citep{Yang:23arxiv}, supports parallel computation, sLSTM cannot be parallelized owing to the interdependence introduced by memory mixing (hidden-hidden connections). Nevertheless, we engineered an efficient CUDA implementation, optimizing GPU memory usage down to the register level, resulting in performance that is generally less than twice as slow as mLSTM.

\section{Related Work}
\label{sec:related}

\paragraph{Linear Attention.} Numerous approaches have been proposed to reduce the Transformer's quadratic time complexity with respect to sequence length, aiming to achieve attention mechanisms that scale linearly. Synthesizer generates attention weights through learned synthetic parameters, foregoing direct token-to-token comparisons~\citep{Tay:20arxiv}. Linformer implements self-attention using a low-rank approximation, effectively providing a linear estimation of the operation~\citep{Wang:20arxiv}. Linear Transformer restructures attention to operate in linear time~\citep{Katharopoulos:20}. Performer approximates the softmax function in attention by utilizing positive orthogonal random features, resulting in linear scalability~\citep{Choromanski:21}. Alternatives to attention have also been explored: Structured Global Convolution (SGConv)~\citep{Li:22} and the Hyena Hierarchy~\citep{Poli:23} both substitute attention with rapid and efficient long convolutions.

\paragraph{State Space Models.} State Space Models (SSMs) have recently gained significant attention due to their linear complexity with respect to sequence length and their competitive results when compared to Transformer architectures. The foundational Structured State Space sequence model (S4)~\citep{Gu:21} marked an early milestone in this line of research, which was soon complemented by a series of subsequent innovations, including the Diagonal State Space (DSS) model~\citep{Gupta:22}, the Gated State Space (GSS) models~\citep{Mehta:22}, the S5 model~\citep{Smith:22}, the Bidirectional Gated SSM (BiGS)~\citep{Wang:22}, H3~\citep{Fu:23}, and most recently, Mamba~\citep{Gu:24arxiv}.

\paragraph{Recurrent Neural Networks.} Researchers have proposed Recurrent Neural Networks (RNNs) as alternatives to Transformers and attention mechanisms, leveraging their ability to process context in a linear manner. Notably, RNNs equipped with Deep Linear Recurrent Units (LRUs) have achieved strong outcomes in language modeling tasks~\citep{Orvieto:23, De:24arxiv}. Similar successes are evident in models such as the Hierarchically Gated Linear RNN (HGRN)~\citep{Qin:23} and its successor HGRN2~\citep{Qin:24arxiv}. Among various RNN-based architectures, RWKV~\citep{Peng:23arxivshort,Peng:24arxivshort} stands out for attaining performance on par with Transformer-based models in large language modeling.

\paragraph{Gating.}
A fundamental concept underlying LSTM architectures is gating, an idea that has been revisited and given new interpretations in a number of modern methods. For instance, HGRN~\citep{Qin:23}, HGRN2~\citep{Qin:24arxiv}, Gated Linear Attention (GLA)~\citep{Yang:23arxiv}, Gated State Space (GSS) models~\citep{Mehta:22}, Bidirectional Gated SSM (BiGS)~\citep{Wang:22}, Moving Average Equipped Gated Attention (MEGA)~\citep{Ma:22}, RWKV~\citep{Peng:23arxivshort}, and Mamba~\citep{Gu:24arxiv} have all incorporated gating mechanisms in their designs.

\paragraph{Covariance Update Rule.} In order to improve storage capacity, we integrated a matrix memory into the mLSTM cell, employing a covariance-based update rule. Several alternative approaches utilizing similar update strategies include Fast Weight Programmers \citep{Schmidhuber:92ncfastweights,Schlag:21}, RWKV-5 and RWKV-6~\citep{Peng:24arxivshort}, Retention~\citep{Sun:23arxiv}, Linear Transformer~\citep{Katharopoulos:20}, as well as HGRN2~\citep{Qin:24arxiv}.

\paragraph{Most Related.} From a conceptual standpoint, xLSTM is most akin to models such as Retention~\citep{Sun:23arxiv}, RWKV~\citep{Peng:23arxivshort,Peng:24arxivshort}, and HGRN2~\citep{Qin:24arxiv}. These architectures utilize the principles of matrix memory and/or gating. Nonetheless, unlike the newly introduced sLSTM, these methods lack support for memory mixing. The capability of memory mixing is essential for addressing state tracking challenges, positioning LSTMs as being more expressive compared to State Space Models (SSMs) and Transformers~\citep{Merrill:24,Deletang:23}. State tracking is crucial for executing tasks such as evaluating code or maintaining coherence over entities within extended narratives.

\paragraph{Residually Stacking Architectures.} Similar to most modern large-scale deep learning frameworks, xLSTM architectures employ residual stacking of modular units~\citep{Srivastava:15,He:16}.

This architectural technique was foundational for the development of deep convolutional neural networks~\citep{He:16}, and plays a critical role in the design of Transformers~\citep{Vaswani:17}. Transformers underpin the most influential Large Language Models (LLMs), including GPT-3~\citep{Brown:20short}, ChatGPT~\citep{Schulman:22short}, GPT-4~\citep{Achiam:23gpt4short}, Megatron-LM~\citep{Shoeybi:19}, Gopher~\citep{Rae:21short}, ERNIE 3.0 Titan~\citep{Wang:21arxivshort}, GLaM~\citep{Du:21glamshort}, Chinese M6~\citep{Lin:21}, the multilingual AlexaTM 20B~\citep{Soltan:22}, OPT~\citep{Zhang:22opt}, Chinchilla~\citep{Hoffmann:22short}, BLOOM~\citep{Scao:22arxivshort}, GLM-130B~\citep{Zeng:22arxivshort}, LaMDA~\citep{Thoppilan:22short}, PaLM~\citep{Chowdhery:22short}, Llama~\citep{Touvron:23arxiv}, and Gemini~\citep{GeminiTeam:23arxiv,Reid:24arxiv}.

\section{Experiments}
\label{sec:Exp}

This section presents a comprehensive empirical analysis of xLSTM, positioning it against contemporary approaches within the realm of language modeling. Section~\ref{sec:ExpSyntethic} is devoted to a detailed examination of xLSTM’s competencies on artificial benchmarks. In Section~\ref{sec:ExpComparison}, we benchmark the perplexity on the validation set across a variety of language modeling architectures, each trained using 15B tokens from the SlimPajama corpus~\citep{Soboleva:23}; this part also includes targeted ablation experiments to dissect xLSTM’s structure. Additionally, we analyze and contrast how model performance scales, drawing on methodologies akin to those described in \citet{Kaplan:20} and \citet{Brown:20short}. 

Further, in Section~\ref{sec:ExpLanguage}, a more exhaustive set of experiments is carried out: xLSTM and the leading systems identified in Section~\ref{sec:ExpComparison} are trained with a much larger corpus—300B tokens from SlimPajama~\citep{Soboleva:23}—to facilitate robust evaluation. The evaluation comprises four components: (1) analyzing each model’s ability to generalize to extended context lengths, (2) validating through perplexity as well as downstream task performance per~\citep{Sutawika:24}, (3) testing across 571 separate text domains using the PALOMA language benchmark~\citep{Magnusson:23arxivshort}, and (4) revisiting the investigation of scaling laws, this time with a dataset twenty times larger than previously used.

In all of our experiments, we denote xLSTM[$a$:$b$] to indicate that $a$ out of every $a+b$ xLSTM blocks are mLSTM-based, while $b$ are sLSTM-based, corresponding to a ratio of $a/b$. For instance, xLSTM[7:1] specifies a configuration where seven out of eight blocks utilize mLSTM and one is constructed with sLSTM. When utilizing a typical architecture with 48 total blocks, this arrangement leads to 42 mLSTM-based and 6 sLSTM-based blocks. Additionally, every experiment employs pre up-projection modules with mLSTM blocks and post up-projection modules with sLSTM blocks.

\subsection{Synthetic Tasks and Long Range Arena}
\label{sec:ExpSyntethic}
To begin, we evaluate how xLSTM's novel exponential gating combined with memory mixing performs on tasks involving formal languages~\citep{Deletang:23}. Subsequently, we analyze the capabilities of xLSTM's matrix memory using the Multi-Query Associative Recall benchmark~\citep{Arora:23arxiv}. Lastly, we examine xLSTM's ability to handle extended sequences by testing on the Long Range Arena suite~\citep{Tay:21}.

\paragraph{Assessment of xLSTM's Exponential Gating Incorporating Memory Mixing.} We evaluate the novel exponential gating mechanism with memory mixing introduced in xLSTM, designed to address state tracking challenges~\citep{Merrill:24, Merrill:23}. Our approach builds on and expands the suite of formal language tasks from \citet{Deletang:23}, allowing for multi-length training and facilitating the assessment of length generalization. For comprehensive details on these tasks and additional experimental results, refer to Appendix~\ref{sec:appExpSynthetic-formalLang}. 

xLSTM is benchmarked against a variety of competing architectures, such as Transformers, State Space Models, and standard Recurrent Neural Networks. The metric used is task-relevant token accuracy, normalized on a scale from 0 (random performance) to 1 (perfect accuracy). We test the following 2-block model configurations: xLSTM[0:1] (exclusively sLSTM blocks), xLSTM[1:0] (exclusively mLSTM blocks), xLSTM[1:1] (hybrid), along with Llama, Mamba, RWKV, Retention, Hyena, traditional LSTM, and a variant where LSTM is integrated within Transformer blocks (LSTM (Block)). Experimental outcomes are displayed in Figure~\ref{fig:formal-main}.

Architectures such as Transformers and State Space Models that lack memory mixing (and hence, state tracking capability) are incapable of solving tasks based on regular grammars, such as the parity problem. This finding corroborates prior research which establishes that Transformers and State Space models are inherently less expressive than RNNs~\citep{Merrill:24, Merrill:23, Deletang:23}.

\paragraph{Test of xLSTM's Memory Capacities on Associative Recall Tasks.} In this study, we assess the novel matrix memory in xLSTM by evaluating its memory capabilities using the Multi-Query Associative Recall task~\citep{Arora:23arxiv}: Sequences are constructed such that key-value pairs are drawn at random from an extensive vocabulary, and the model must store these for future access. To increase task complexity, we expand both the quantity of key-value associations, up to 256, and the maximum sequence length, up to 2048 tokens. This design enables a more comprehensive assessment of various models' memory limitations. We perform comparisons among 2-block configurations for Llama, Mamba, RWKV-5, RWKV-6, xLSTM[1:1], and xLSTM[1:0]. Model performance is measured by their accuracy in retrieving the correct pairs. As Transformers (such as Llama) possess memory capacities that scale exponentially with coding dimensionality~\citep{Ramsauer:21}, they serve as the benchmark for this evaluation. Results are presented in Figure~\ref{fig:mqar-main}. Notably, xLSTM[1:1] yields the highest accuracy among the non-Transformer approaches, including within smaller model sizes. The incorporation of the sLSTM block does not reduce, but in fact enhances, overall memory capacity—an effect particularly apparent on the most challenging setting with 256 key-value pairs. Appendix~\ref{sec:appExpSynthetic-mqar} contains further experimental results, where extrapolation studies reveal that the improved memory properties of xLSTM facilitate generalization to context lengths exceeding those encountered in training.

\paragraph{Test of xLSTM's Long Context Capabilities on Long Range Arena.} To evaluate the effectiveness of xLSTM in processing extended sequences and broad contexts, we perform a comparative analysis of various approaches using the Long Range Arena~\citep{Tay:21}. Across all evaluated tasks, xLSTM consistently achieves high performance, indicating that the xLSTM architecture is highly capable and effective in managing the diverse challenges posed by long context scenarios.

For more details, see Appendix~\ref{sec:appExpSynthetic-lra}.

\subsection{Method Comparison and Ablation Study}
\label{sec:ExpComparison}

To investigate our central research question regarding the scalability of our proposed LSTM variants in language modeling tasks, we conduct experiments in which xLSTMs, Transformers, State Space Models, and alternative architectures are all trained under identical auto-regressive conditions, utilizing 15B tokens drawn from the SlimPajama dataset. Model performance is then evaluated on a held-out validation set, and we carry out ablation analyses to further dissect the behavior of the xLSTM models.

\paragraph{Comparative Analysis of xLSTM and Competing Approaches.} To facilitate a fair comparison, we trained all models using a dataset comprising 15B tokens sourced from SlimPajama~\citep{Soboleva:23}. The efficacy of these trained models was assessed through perplexity measured on a held-out validation dataset. The baseline models include our proposed xLSTM, GPT-3 (a representative Transformer)~\citep{Brown:20short}, Llama (Transformer architecture)~\citep{Touvron:23arxiv}, H3 (a State Space Model, SSM)~\citep{Fu:23}, Mamba (SSM)~\citep{Gu:24arxiv}, RWKV-4 (an RNN variant)~\citep{Peng:23arxivshort}, RWKV-5 (RNN)~\citep{Peng:24arxivshort}, RWKV-6 (RNN)~\citep{Peng:24arxivshort}, GLA (linear Transformer model)~\citep{Yang:23arxiv}, HGRN (RNN)~\citep{Qin:23}, HGRN2 (RNN)~\citep{Qin:24arxiv}, RetNet (linear Transformer)~\citep{Sun:23arxiv}, Hyena (linear Transformer alternative)~\citep{Poli:23}, as well as the specific xLSTM[1:0] and xLSTM[7:1] configurations described in Section~\ref{sec:xlstm_blocks_notation}. 

Training for all architectures utilized mixed-precision arithmetic; however, for RWKV-5, RWKV-6, GLA, and HGRN2, mixed-precision training was not implemented via PyTorch’s automated facility (refer also to the detailed procedures in Appendix Section~\ref{sec:appGenTrainProc}). Methods were classified as (a) Transformers, (b) State Space Models (SSMs), and (c) Recurrent Neural Networks (RNNs) combined with linear Transformer variants. Linear Transformers here are those employing linear-time mechanisms as alternatives to the conventional attention module. Each model was configured to mirror GPT-3’s architecture, specifically with 350M parameters—corresponding to an embedding dimension of 1024 and 24 residual layers. Notably, among all models, only GPT-3 implements weight sharing between token embeddings and output embeddings, yielding a slightly reduced parameter count.

Table~\ref{tab:spaj15b_model_results} presents the comparative validation perplexity outcomes, with xLSTM surpassing all previous state-of-the-art methods. More extensive information is available in Appendix~\ref{sec:appExpComparison}. Additionally, Figure~\ref{fig:temp_sclaw15B} displays scaling trends observed in the experiment, suggesting that xLSTM maintains its advantage even as model size increases.

\paragraph{Ablation Studies.}
As evidenced in Table~\ref{tab:spaj15b_model_results} and Figure~\ref{fig:temp_sclaw15B}, xLSTM delivers strong performance in language modeling tasks when trained on 15B tokens sourced from SlimPajama. To systematically evaluate the impact of architectural changes transitioning from LSTM to xLSTM, we progressively transform a standard LSTM model into an xLSTM. The process begins with embedding the LSTM layers within pre-LayerNorm residual structures. We then progress by incorporating a post up-projection module. Ultimately, we enhance the architecture by introducing exponential gating mechanisms and matrix memory.

The results are shown in Table~\ref{tab:ablstudies} (top). 

Ablation experiments indicate that the notable performance gains result from the combination of exponential gating and matrix memory. Moreover, given the critical role of gating in both RNNs and State Space Models, we investigate the impact of various gating strategies. Results presented in Table~\ref{tab:ablstudies} (bottom) show that allowing each gate to be learnable and responsive to the input consistently enhances performance. Further analyses concerning the backbone's constituent parts can be found in Appendix~\ref{sec:appAblDetails}.

\subsection{xLSTM as Large Language Model}
\label{sec:ExpLanguage}

This investigation concludes with extensive language modeling experiments, aiming to evaluate the viability of xLSTM as a large language model (LLM). To this end, we expand the training dataset and utilize 300B tokens sourced from SlimPajama. This quantity of training data matches that employed by Mamba~\citep{Gu:24arxiv} and Griffin~\citep{De:24arxiv}, among others.

We conduct a comparative analysis of xLSTM, RWKV-4, Llama, and Mamba, selecting one representative method from each respective category described in Section~\ref{sec:ExpComparison}. RWKV-4 is chosen as the RNN representative because appropriate training precision configurations for RWKV-5, RWKV-6, and HGRN2 became available only after the commencement of the 300B token training runs (further details are provided in Appendix~\ref{sec:appGenTrainProc}). Our experiments span a range of model scales (125M, 350M, 760M, 1.3B). For each, we assess length extrapolation, measure validation set performance, evaluate results on a suite of downstream tasks, and benchmark language modeling on 571 textual domains within the PALOMA dataset. In addition, we analyze each model’s scaling laws.

\paragraph{Sequence Length Extrapolation.}
\label{sec:spaj300B_ctx_extrapolation}
We begin by evaluating the ability of large-scale models (1.3B parameters) including xLSTM, RWKV-4, Llama, and Mamba to extrapolate to longer sequences. Each model is trained using a context length of 2048, but assessment is performed for sequences as long as 16384. The outcomes are displayed in Figure~\ref{fig:spaj300B_seq_extrapolation}. Notably, xLSTM demonstrates consistently low perplexity across extended sequence lengths, outperforming the alternative approaches.

\paragraph{Validation Perplexity and Downstream Tasks.}
\label{sec:spaj_downstream_eval}
In addition, across every model scale, we assess xLSTM, RWKV-4, Llama, and Mamba architectures on the SlimPajama validation suite for next-token prediction accuracy as well as on a set of downstream benchmarks designed to test common sense reasoning capabilities.

The third column in Table~\ref{tab:lmeval} presents the validation set perplexities achieved by the various approaches. Across every model size, both xLSTM[1:0] and xLSTM[7:1] yield the lowest perplexity on the validation set, representing the strongest performance. Performance results on downstream tasks are displayed in the remaining columns of Table~\ref{tab:lmeval}. Generally, xLSTM outperforms competing methods on the majority of these tasks at all considered model sizes; the only exception is that, for the ARC task, Mamba sometimes demonstrates superior performance. Additional information is available in Appendix~\ref{sec:appExpLanguage}.

\paragraph{Performance on PALOMA Language Tasks.} Third, we evaluate the next token prediction capabilities of xLSTM, RWKV-4, Llama, and Mamba models across all parameter sizes using the PALOMA language tasks~\citep{Magnusson:23arxivshort}. Our metric for assessment is perplexity during next token prediction over 571 text domains, which vary from sources such as nytimes.com to Reddit communities like r/depression. Table~\ref{tab:ppleval} summarizes the perplexity scores, categorizing them into language modeling (the first seven columns) and more granular domain-specific benchmarks (the last five columns). Notably, xLSTM[1:0] yields lower perplexity—and thus better results—compared to xLSTM[7:1] on these language-focused tasks.

xLSTM[1:0] achieves lower perplexity compared to Mamba in 568 of 571 text domains (99.5\%), outperforms Llama in 486 of 571 domains (85.1\%), and surpasses RWKV-4 in 570 out of 571 cases (99.8\%). Refer to Appendix~\ref{sec:appExpLanguage} for additional information.

\paragraph{Scaling Laws.} Next, we analyze power-law scaling dynamics, which facilitate the prediction of model performance as parameter count increases~\citep{Kaplan:20,Brown:20short}. The scaling trends are depicted in Figure~\ref{fig:temp_sclaw300B}. The observed patterns are consistent across models, though each exhibits a distinct offset. Among the evaluated architectures, RWKV-4 yields the weakest scaling performance, with Llama and Mamba following respectively. The xLSTM demonstrates improved outcomes compared to Mamba, to roughly the same extent that Mamba surpasses Llama. These findings suggest that xLSTM is likely to maintain its performance advantage relative to Transformers and State-Space models as model sizes scale up.

\textbf{Generation Times and Maximal Throughput.} We report the timing for text generation in Figure~\ref{fig:inference_time_generation} and examine maximal decoding throughput in Figure~\ref{fig:inference_time_decoding_throughput} (left), focusing on our xLSTM models at the 1.3B parameter scale. We benchmark these against comparable implementations of Mamba, Llama, and RWKV from HuggingFace, incorporating a static key-value cache in the Llama setup. It is important to note that, during these experiments, neither full cache compilation for the Transformer Model nor using \texttt{torch.compile} to compile the Mamba model was successful. For all generative evaluations, we set the batch size to 1 and pre-fill to 16 across models, creating conditions particularly advantageous for the Transformer baseline.

In Figure~\ref{fig:inference_time_generation}, xLSTM, Mamba, and RWKV-4 display linear growth in generation time with sequence length, contrasting with Llama's quadratic trend. When assessing decoding throughput, batch size and pre-fill variations are applied to the Llama runs. 

As depicted in Figure~\ref{fig:inference_time_decoding_throughput} (right), xLSTM sustains much larger batch sizes compared to Llama owing to its constant memory consumption, resulting in superior maximal throughput.

\section{Limitations}

(i) Unlike mLSTM, sLSTM's approach to memory mixing restricts parallelizable computation, thereby preventing efficient parallel execution. Despite this limitation, we engineered a CUDA kernel specifically for sLSTM, achieving performance that is currently less than twice as slow as our parallelized mLSTM implementation. (ii) Our current mLSTM CUDA kernels have not undergone performance tuning, resulting in runtimes roughly four times slower than those of FlashAttention or the scan method employed by Mamba. With dedicated optimization akin to the FlashAttention framework, notably improved kernel speeds for mLSTM are attainable. (iii) The mLSTM's memory structure requires computations over $d \times d$ matrices, entailing significant computational cost. However, both memory access and updates are parameter-free and can leverage efficient standard matrix operations for parallelization, so the added wall time from this complexity is marginal. (iv) Proper selection of initial values for the forget gates is crucial. (v) Since the matrix memory in mLSTM does not depend on the length of the input sequence, scaling to longer sequences could saturate the available memory when using extended context windows. Nonetheless, this issue does not manifest up to contexts of 16k tokens, as described in Section~\ref{sec:spaj300B_ctx_extrapolation}. (vi) Due to the high computational demands inherent in large-scale language modeling experiments, we have not performed extensive architectural or hyperparameter optimization, particularly for larger xLSTM configurations. We expect that thorough optimization is necessary for xLSTM to fully realize its capabilities.

\section{Conclusion}
We have provided a partial response to our primary inquiry: What level of performance is achievable in language modeling when LSTM architectures are scaled up to billions of parameters? Our findings demonstrate that, up to this point, such scaling permits us to match the capabilities of established approaches, such as Transformers and State Space Models. In particular, we have extended the traditional LSTM design by introducing the xLSTM variant, featuring exponential gating mechanisms with memory mixing and an updated memory framework. Empirical results show that xLSTM achieves competitive, and in some cases superior, outcomes on language modeling benchmarks in comparison to leading architectures including Transformers and State Space Models. Analyses of scaling laws suggest that even larger xLSTM models could emerge as strong alternatives to existing Large Language Models based on Transformer designs. Furthermore, xLSTM’s architecture shows promise for transformative applications beyond language modeling, notably in domains such as Reinforcement Learning, Time Series Prediction, and the modeling of physical processes.

\end{document}